\documentclass[12pt, ukrainian]{article}
\usepackage[a4paper]{geometry}
\setlength\parindent{0mm}
\setlength\parskip{4mm}
\geometry{top=1.5cm,bottom=1.5cm,left=1.3cm,right=1.5cm}
\pagestyle{empty}
\usepackage[T2A]{fontenc}
\usepackage[utf8]{inputenc}
\usepackage[ukrainian]{babel}
\usepackage{graphicx}
\usepackage{caption}
\usepackage{caption}
\DeclareCaptionFormat{nocaption}{}
\captionsetup{format=nocaption,aboveskip=0pt,belowskip=0pt}
\usepackage[Export]{adjustbox}
\adjustboxset{max size={0.9\linewidth}{0.9\paperheight}}
\usepackage{verbatim, listings, clrscode3e, fancyvrb}
\def\code#1{\Verb!#1!}
\begin{document}
\begin {large}

\textbf{01.12.2023 Контрольна робота, варіант  \No { 30 }}\par


{
{\begin {center}\textbf{Завдання 1}\end {center}}\par
По яких днях тижня відбувались практичні з предмету?\par
{\begin {center}\textbf{Завдання 2}\end {center}}\par
Як зовуть (ім'я, по батькові) викладача, що вів практичні заняття?\par
{\begin {center}\textbf{Завдання 3}\end {center}}\par
Виберіть всі правильні твердження, якщо такі тут є.\par
\vspace{2mm}
( 1 ) Послідовність може мати дві різних експоненційних твірні.%bad

( 2 ) Питання, чи існує послідовність, що має дві різні твірні, є відкритим.%bad

( 3 ) Паросполучення можна промоделювати сполукою з повтореннями з обмеженнями на розташування елементів.%bad

( 4 ) Задачу знаходження кількості способів сформувати новорічні подарункові набори з $n$ цукерок, що складаються з цукерок 13 видів, де цукерок кожного вигляду не більше 20, можна розв’язати методом рекурентних співвідношень.%good

( 5 ) Турові поліноми можна розглядати не тільки для квадратних дошок.%good

( 6 ) Послідовність, задана рекурентним співвідношенням $$a_{n+1}=a_{n}a_{0}+a_{n-1}a_{1}+\ldots+a_0a_{n}$$ може не бути послідовністю чисел Каталана.%good

( 7 ) Коли кажуть про твірну функцію послідовності, то на увазі завжди мають експоненційну твірну.%bad

( 8 ) Задачу знаходження кількості способів складання кодових слів довжиною $n$, що складаються з десяткових цифр, де парні цифри зустрічаються непарну кількість разів, а парні цифри — непарну кількість разів, можна вирішити за допомогою твірної.%bad

( 9 ) Паросполучення можна промоделювати розміщенням з обмеженнями на розташування елементів.%good

( 10 ) Послідовність, задана рекурентним співвідношенням $$a_{n+1}=a_{n}a_{0}+a_{n-1}a_{1}+\ldots+a_0a_{n}$$ є послідовністю чисел Каталана.%bad\par
}
\end {large}

\vspace{4mm}
%end
\newpage
\end{document}

